\chapter{EVALUASI}

\petunjuk{Judul bab ini dapat disesuaikan dengan karakter tugas akhir. Jika evaluasi telah digabungkan dalam BAB IV, bab ini dapat diubah menjadi bab lain yang relevan atau dihapus dengan tetap menyesuaikan penomoran.}

\section{Evaluasi terhadap Tujuan}

Evaluasi menjelaskan sejauh mana setiap tujuan penelitian telah dicapai. Gunakan indikator yang telah didefinisikan pada metodologi dan hindari memperkenalkan metrik baru tanpa penjelasan.

\section{Keterbatasan}

Tuliskan keterbatasan data, metode, perangkat, ruang lingkup, dan proses pengujian secara jujur. Keterbatasan bukan sekadar daftar kekurangan, tetapi menjadi dasar untuk memahami ruang berlaku hasil penelitian.

\section{Implikasi}

Jelaskan implikasi hasil bagi pengguna, organisasi, pengembang, atau penelitian lanjutan sesuai dengan konteks tugas akhir.
